% Source: https://artofproblemsolving.com/community/c3h1429601_simultaneous_equations
% Stable identity: aops-topic-1429601-q01
Set the right sides equal. Since $(x+8)\left(x^2+2\right)=x^3+8x^2+2x+16$,
\[
4(x+2)y=x^3+3x^2+18x+32=(x+2)\left(x^2+x+16\right).
\]
\textbf{Case $x=-2$.} The first equation gives $y^2=6\cdot6=36$, so $y=\pm6$, and the second gives $y^2=0+20+32-16=36$ as well. So $(-2,6)$ and $(-2,-6)$ are solutions.

\textbf{Case $x\ne-2$.} Then $y=\dfrac{x^2+x+16}{4}$, and the first equation becomes
\[
\left(x^2+x+16\right)^2=16(x+8)\left(x^2+2\right),
\]
\[
x^4+2x^3+33x^2+32x+256=16x^3+128x^2+32x+256,
\]
\[
x^4-14x^3-95x^2=0\ \Longrightarrow\ x^2(x-19)(x+5)=0 .
\]
So $x=0$, $19$ or $-5$, with $y=\dfrac{x^2+x+16}{4}=4$, $99$, $9$ respectively. Each pair satisfies both equations (the second equation was used only through the difference, and the first holds by construction).

The solutions are
\[
\boxed{(x,y)=(-2,6),\ (-2,-6),\ (0,4),\ (19,99),\ (-5,9)} .
\]
